\subsection{IA et cinéma}

\begin{itemize}
	\item Archives audiovisuelles  
	
	"When you start looking at moving images and doing ML on them, your object of study
	is no longer an entire film but shots or scenes or still images. And that suddenly brings
	in important questions for archives and libraries. When the record itself changes…that
	raises huge questions that folks are not grappling with. Not just enriching a particular
	object but changing the object". 
	\footcite{cordell}
	
	Difficultés de ce type de contenu : 
	"moving image data, which is a particularly difficult data type to atomize."
	\footcite{cordell}
	
	\item Cas d'usage 
	
	Exemple de projet d'extraction d'entités depuis des images en mouvement : \footcite{arnold2019}
	méthode du distant viewing: reconnaissance des visages, classification des plans. Le montage et le cadrage servent à costruire la place des personnages dans le récit. 
	"« One of the first tasks in understanding a video file is to break up the sequence of frames into distinct units known as shots ».
	« We first split an input video file into individual shots, then locate faces within every frame, identify the faces in a frame, and finally put all of the information together to cluster shots based on camera angles and distances ».\footcite{arnold2019} 
	Shot boundary detection : détection automatique des coupures entre les plans dans une vidéo 
	Shot semantics : description de ce qui est présent dans l'image et à quel endroit de l'image et ainsi relier l'image à des catégories de plan (shot type).
	
	problématique principale pour indexer : il faut déjà séparer la vidéo en tous les plans qui la composent pour en extraire les éléments : 
	"One of the first tasks in understanding a video file is to break up the sequence of frames into distinct units known as shots—or uninterrupted footage between
	cuts—in a process known as shot boundary detection."\footcite{arnold2019}
	
	exemple de pipeline pour détecter les coupures entre les plans : 
	"the vast majority of cuts in most film and television
	episodes consist of abrupt transitions between alternative camera shots. These are both easy to define and relatively easy to detect. The algorithm we use in our analysis uses a combination of histogram differences and down sampled pixel intensities. Our method functions by measuring the extent to which two successive shots are different from one another in terms of the general color palette and the distribution of brightness over the frame. If these differences fall above a specified threshold, we indicate the presence of a cut.22 We also added special logic to avoid erroneous cuts during fade-in and fade-out events." \footcite{arnold2019}
	
	face detection models : Faster R-CNN and FAREC-CNN \footcite{arnold2019}
	
	shot classification : il est possible de classifier le type de plan à l'image en utilisant le placement des visages : "In this final step we no longer work directly with the visual materials and build a meta-algorithm that takes detected shot breaks and faces as input and outputs categorical predictions for each shot". \footcite{arnold2019}
	Il s'agit de définir tous les types de plans età partir de ces définitions, construire un jeu de donnes qui sera testé : close-up, over the shoulder...
	
	Audiovisuel Processing (AVP) : ensemble de technologies appliquées à l'analyse de contenus vidéos et audio (comprend computer vision et ASR automatic speech recognition) 
	\footcite{noord2021}
	
	constat : malgré de grandes avancées dans le processing audiovisuel, ces outils restent largement inaccessibles et ne sont pas donc pas adoptés dans le milieu des archives : « Despite the massive progress made in the performance levels and overall availability of Audiovisual Processing (AVP) technologies, such as Computer Vision (CV) and Automatic Speech Recognition (ASR), the availability of these tools for heritage institutions is still lagging behind. »
	\footcite{noord2021}
	
	ImageJ : open source visual analysis software suite 
	\footcite{noord2021}
	Detectron2 : détection de poins corporels dans l'image \footcite{noord2021}
	
	les AVP ne sont pas accessible en libre accès aux chercheurs en humanité numériques. Les algorithmes ont besoin d'un accès direct aux données, or, pour des raisons de sécurité des données et des droits de propriétés intelectuelles auxquels sont soumis les archives, les données des archives ne peuven tpas quitter leur environnement numérique. 
	donc : "A side effect of this is that for the analysis of AV material commercially available tools cannot be used, as these are primarily cloud based. To tackle this challenge it is necessary that any AVP infrastructure can
	be integrated with a local archival infrastructure, i.e.,
	we need to bring the algorithms to the data, rather than
	the data to the algorithms"
	\footcite{noord2021}
	
	au lieu de transférer les fichiers vers un serviec externe, l'algorithme doit être déployé localement au sien même de l'ifrastructure de l'archive. 
	
	exemple de l'architecture DANe au sien du projet CLARIAH : l'architecture est divisée en plusieurs "workers" et une file d'attete. si l'u des worers s'arrêtent de fonctionner, les autres continuent leur tâche. 
	Cette architecture produit deux types de données : des métadonnées structurées directement epxloitables, et des vecturs de caractéristiques (= représentations abstraites des données utilisables pour des requêtes de similarité : ouvre des possibilité de recherche par ressemblance son et vidéo.)
	\footcite{noord2021}
	
	"there is a near-linear relationship between quality of transcriptions and search performance.» : importance de la qualité de la transcription audiovisuelle pour améliorer la recherche.
	\footcite{noord2021}
	
	"the application of AVP tools to large image datasets significantly scales up the level of analysis from ‘close’ to ‘distant’ viewing, allowing scholars to observe patterns in the data that are invisible in the traditional, interpretative analyses of small sample sets.» 
	le distant viewing permet de passer d'une lecture individuelle d'items à une analyse d'éléments à l'échelle d'un corpus. 
	\footcite{noord2021}
	
	"« Digital tools for paper resources have tended not to consider the strong presence of visual elements in them. Due to the lack of practically available AVP tools to digital humanities, scholarly research has ‘grossly neglected visual content, causing a lopsided representation of all sorts of digitized archives’. »
	Sous représentation des éléments visuels dans l'indexation, même pour une collection audidiovisuelle. Les métadonnées textuelles liéees aux films sont nsuffisantes pour enrichir leurs métadonnées : titre, synopsis...
	\footcite{noord2021}
	
	
	\item Indexation des archives audiovisuelles 
	
	Premier outil d'indexation des images : Perzeptron de Franz Rozenblatt en 1957 : machine de reconnaissance d'images.
	
	Pour l'indexation des archives audiovisuelles : apprentissage supervisé, non supervisé ou en profondeur : sont appliqués ou combinés. 
	
	"Dans le cas des archives audiovisuelles, on voit en outre très clairement	comment les besoins des utilisateur·rice·s sont pris en compte lors de la conception de	l'indexation automatique en profondeur et comment les données et les traces des	utilisateur·rice·s sont utilisées pour optimiser l'accès aux archives audiovisuelles (mot-clé :	système de recommandation)."
	\footcite{zotero-item-584}
	
	"L'indexation, la recherche et le filtrage basés sur le texte de	contenus audiovisuels pertinents se heurtent par nature à des limites : Les contenus	audiovisuels doivent également pouvoir être recherchés à l'aide de méthodes audiovisuelles	(par exemple la navigation par images, Query by Example)."
	\footcite{zotero-item-584}
	
	"L'identification de contenus simples d'images ou de sons basée sur des descripteurs
	techniques ou des attributs dérivés a tendance à être de plus en plus laissée à la machine
	apprenante (par exemple « photo de groupe », « personne avec un casque sur la tête »,
	termes parlés), tandis que l'identification plus poussée de personnes, de corps, de lieux,
	d'événements ou de sujets concrets est encore très sujette à erreurs."
	\footcite{zotero-item-584}
	
	"When it comes to processing
	image-based records in archives archives, supervised learning models can be	helpful for identification and classification, while unsupervised learning	models are best suited for processing large amounts of data for creating	generic descriptions, and highlighting possible connections between records."
	\footcite{fewster}
	
	\item Outils 
	
	YOLO (you only look once)
	Algorithme qui utilise CNNs pour entourer des objets dans des images et les classer dans des catégories 
	il s'agit d'un single shot object detector : il se base sur une seule analyse de l'image dans sa globalité 
	facial recognition, scene analysis, object tracking \footcite{fewster}
	
	two shot objects detector : multiple image reviews 
	Exemple : Mask	Region-based CNN (R-CNN) algorithm. "uses deep neural networks to	identify regions of interest in the photograph and build boxes around objects	found throughout the image" \footcite{fewster}
	Cette manière de faire est appelée "semantic segmentation" : "identifies semantic objects	without specifying how many instances of that object are present in the	image" \footcite{fewster}
	Par exemple, s'il y a 5 personnes dans une image, la segmentation sémantique les identifie comme un même objet sémantique dans l'image.
	
	Instance segmentation : "identifies all the semantic objects in	an image and segments them into specific instances to accurately distinguish	between each object of a similar semantic category" (c'est ce que fait Mask RCNN)\footcite{fewster}
	
	Computer–Aided Metadata generation for Photo Archives Initiative (CAMPI) \footcite{fewster} production de tags sur des images et recommandation d'images similaires 
	
	
	\item Biais et limites 
	
	"For instance, racial bias has been an ongoing challenge with	computer vision tools, particularly with facial recognition software. A study	published by MIT Media Lab in 2018 found that facial recognition software	had an average error rate of 0.8 percent for light-skinned men, versus a 34.7 percent	error rate for darker-skinned women" \footcite{fewster}
	"societal biases are often engrained in	AI model training datasets, which lead to the models reproducing those	biases in their outputs" \footcite{fewster}
	"In addition to this, the training data of computer	vision models has relied mostly on visual content available to be harvested	from the Internet, which in its majority corresponds to content developed	during the last three decades. These models would underperform on images	that are of a more historical character, commonly found in visual archives,	since they are underrepresented in the training data" \footcite{fewster}
	
	IL est important pour les archivistes de savoir comment ont été entraînés les modèles qu'ils utilisent, sur quelles données.
	
	la vidéo : souvent de moins bonne qualité que l'image fixe, donc les résultats d'indexation automatique peuvent être moins bons
	
	contrairement aux documents textuels qui comportent souvent des éléments d'indexation dans leur forme (couverture de livre, ) les archives audiovisuelles n'en contiennent peu ou pas, ce qui les rend opaques et peu réceptives aux pratiques d'indexation. Il s'agit alors de mettre en place une « indexation par le contenu ». \footcite{zotero-item-582} 
	
	Une vidéo doit être segmentée avant de pouvoir être indexée : 
	« One of the most common initial steps for video content organisation is segmenting the video into different meaningful parts. Proper segmentation facilitates content organisation, improves accessibility, and enhances the ability to perform analytics on large video datasets. »\footcite{attard2025}
	
	Une idée de pipeline de segmentation : est la détection de scènes avec pyscenedetect (découpage d'une vidéo en segments en détectant les coupes visuelle), puis classification des segments. D'un côté la segmentation syntaxique et de l'autre la classification sémantique. \footcite{attard2025}
	
	le meilleur classifieur d'images avec le plus haut résultat de précision selon attard 2025 est ResNet
 \footcite{attard2025}
 
l'article de attard souligne que  « The study underscores the lack of specialised annotation tools and the high resource demands of sophisticated models, which could present significant obstacles for future research with limited computational resources. » \footcite{attard2025} c'est une limite surtout pour les isntitutions. REjoint l'argument et les limites exposées par \footcite{sullivan2025}

pour justifier la mise en place d'une segmentation : « Automated segmentation accelerates the processing of large volumes of video content and ensures consistency and accuracy in the segmentation process. » \footcite{attard2025} la mise en place d'une automatisation de la segmentation permettrait d'homogénéiser l'indexationdocumentataire : un algorithme segmente les vidéos de la même façon, pas les humains. Mais les vidéos sont toutes différentes : est ce qu'on peut garder la même pratique de segmentation sur des corpus vidéo différents ?

les modèles multimodaux qui combinent audio et vidéo (vivit-ast) ne donnent pas les meilleurs résultats. Il faudrait donc adopter des architectures plus simples, qui demandent moins de calcul d'usage et qui séparent les pipelines. \footcite{attard2025}

absence d'outils d'annotation spécialisés dans la segmentation vidéo. Attard adapte label studio, un outil généraliste. \footcite{attard2025}

le modele ViViT, modèle basé sur des transformers, capture les caractéristiques spatio temporelles de la vidéo, contrairement à Resnet qui analyse chaque image d'un vidéo indépendemment. ViViT extrait "extracts spatio-temporal tokens from the input video, which are then encoded by a series of transformer layers. » \footcite{arnab2021}. C'est un modèle qui tete d'analyser l'espace et le temps dans une vidéo donnée. Vivit adopte donc une approche "video-native". 
Le modèle à choisir dépend de la nature du corpus audiovisuel que l'on souhaite indexer. Si les archives sont des plans fixes qui ne changent que très peu, resnet suffit. 

« ViT has been shown to only be effective when trained on large-scale datasets, as transformers lack some of the inductive biases of convolutional networks, and thus require datasets larger than the common ImageNet ILSRVC dataset to train. » \footcite{arnab2021}. : les transformers n'ont pas les biais des CNN qui leur permettraient de faire des généralisations et de reconnaître un objet où qu'il soit dans l'image. Ils ont donc besoin d'être entrainés sur beaucoup plus de données : c'est une grande limite. Une approhe par calssifieurs d'images plus légers comme ResNet peut donc être préférable. ViViT obtient de très bons résultats mais entraînés sur des corpus vidéos contemporains : vidéos youtube. Diffiilement adaptable à un corpus d'archives historique sauf pré entraînement massif. 

Récemment, une transition architecturale des modèles a produit un basculement vers des modèles basés sur les transformers en vision par ordinateur, remplaçant prograssivement les CNN. Cela explique pourquoi les outils actuellement disponibles comme CLIP, ViViT, AT, sont tous basés sur cette architecture. 

la qualité des images est absolument primordiale et déterminante dans la mise en place d'un pipeline d'indexation automatique : "La question de la qualité des images est
cruciale dans le cadre d’un projet de recherche en vision artificielle, puisque celle-ci a la possiblité d’impacter les performances de l’algorithme utilisé, et d’affecter négativement les interprétations du modèle". \footcite{norindr2023}
	
\end{itemize}